\begin{afterword}
\summaryhead

The post is a dialogue. A narrator asks how evolution, a blind and brutal process, could produce beings capable of love. A second voice, addressed as ``you'', answers each case with an evolutionary explanation: mothers love children who carry their genes, friendships pay off in repeated cooperation, and moral argument grew out of tribal politics. Gandhi is explained as someone who believed and acted on general moral propositions.

The narrator keeps insisting that human kindness and our sense of beauty are surprising. The second voice replies that they seem marvelous only because they are our own values. If there were an objective standard of value outside us that agreed with what evolution happened to produce, that would be too much of a coincidence.

The narrator settles for calling human love a ``moral miracle''. The post ends with a fairy tale in which future children ask where love began, and are told of lovers made by something that did not love.

\responsehead

The post is written as a debate, and the author writes both sides. Every link in either voice goes to one of his own posts. The reader is cast as the voice that wins (``you say''), and the narrator, who feels wonder, is corrected at every turn and finally agrees to a redefinition of his own words. Twice his feeling that something has been lost is diagnosed as a personal failing, and the winning voice reports that it takes more joy in flowers than he does. What the reader is given is a position to occupy: the one who understands, looking at the one who is merely moved.

The argument does not survive its own objector. The narrator says the explanation of Gandhi ``could explain any possible human behavior''. That is correct, and it is never answered. The reply is a hedged list of causes and then a false choice between that list and magic. The central argument, that an objective standard of value matching our evolved values would be ``too much coincidence'', is, in substance, the evolutionary debunking argument Sharon Street published in 2006, two years earlier. It is not credited. It is a live dispute with a literature of replies. The post settles it with rhetorical questions and a standard, moving a ``ghost of perfect emptiness'', that no argument could meet, its own included.

The science is offered with more confidence than it has. The prisoner's dilemma has no single ``solution''. Flowers in general did not evolve by imitating the mating signals of bees; that trick is known only in orchids. The claim that moral argument bears the ``specific design signature'' of the savanna is asserted, not shown. The origin of life is dated to 3.85 billion years, a contested figure, and placed in a tidal pool without the hedge Darwin gave his ``warm little pond''.

The ending does most of the post's work. The question it opened with is answered by a bedtime story. By its author's own account in ``Something to Protect'', such campfire tales satisfy curiosity just as well as true answers do. Inside the story sit two claims the post never argues: our descendants will be ``intelligently designed'' by us, and civilization will outgrow a single galaxy. Inference: this designed future is the ``gift'' of the title, the one word the text never uses.

In short: a debate in which the author writes both sides and makes the objector yield, which borrows an uncredited argument from moral philosophy, gets the natural history of flowers wrong, and ends by delivering its vision of the future as a fairy tale.
\end{afterword}
